\scanpage{2}

\margin{(29.11.83)}

Soit $X$ surf. à bord compacte orientée, à structure torique
\inserted{(i.e. torseur sous $\mathbb U=\{z\in\mathbb C\mid |z|=1\}$)}
donnée sur les comp. connexes du bord.
On définit
\[
\begin{aligned}
\operatorname{SRT}(\unsure{X})={}&
\text{quotient du groupe des autom de }X
\text{ (\deleted{\unsure{en}} resp. l'orientation}\\
&\text{et les structures toriques) par la comp. neutre du}\\
&\text{sous-groupe }\operatorname{Aut}^!(\vec X)
\deleted{\unsure{\mathrm{bijec.}}}
\text{ de Aut qui induit}\\
&\text{l'identité sur }\partial X
\end{aligned}
\]

On a ainsi un \unsure{dévissage} de $\operatorname{SRT}(X)$ en trois étages
\[
\underbrace{S_{\unsure{*}}T^!(X)}_{
\operatorname{Aut}^!(\unsure{\vec X|\overline X})/
\text{comp. neutre}}
\qquad\mathbb U^I\qquad\unsure{\mathfrak S_I}
\qquad I=\pi_0(\partial X)
\]
ou encore
\[
\mathbb Z^I\qquad T(X)\qquad\mathbb U^I\qquad\unsure{\mathfrak S_I}
\]

En fait, on définit can! un sous-groupe \inserted{droit}
$\mathbb R^I$ de $\operatorname{SRT}(X)$
contenant $\mathbb Z^I$, de telle façon qu'on ait
\[
\underbrace{\operatorname{SRT}^!(X)}_{
\substack{\text{noyau de}\\
\operatorname{SRT}^!(\unsure{X|\alpha})\to\unsure{\mathfrak S_I^!|\mathfrak S_I'}}}
=\mathbb R^I.S_*T^!(X)
\]

\diagram{
\begin{tikzcd}[column sep=small,row sep=large]
\mathbb R^I &
\mathord{:}\arrow[l,phantom,"\subset"']\arrow[dr,no head,"\unsure{\lambda}"] &&\\
\mathbb Z^I\arrow[u,"\unsure{\mathrm{can}}"']\arrow[r,hook]&
S_*T^!\arrow[u,hook]\arrow[r,phantom,"\subset"]&
\operatorname{SRT}^!\arrow[r,phantom,"\subset"]&
\operatorname{SRT}
\end{tikzcd}
}
\margin{\deleted{\unsure{et propre}}}

i.e. on a

\margin{\uline{NB} $\operatorname{SUT}(X)\simeq
T(X)\times\underset{\operatorname{Aut}((T_i)_{i\in I})}{\unsure{\mathfrak S_I}}$}

\margin{\unsure{jusqu'ici on ne fait intervenir}
que $S_*T^!(X)$; à l'exclusion de $S_*T(X)$,
dont la défini. exige une str. suppl.
$\uline{s}$, cf ci-dessous}

\diagram{
\begin{tikzcd}[column sep=small,row sep=large]
&\operatorname{SRT}(X)\arrow[d,no head,"\unsure{\mathfrak S_I}"]&\\
&\mathbb R^I.\unsure{S_*T^!}\arrow[dl,no head,"T^!"']
 \arrow[dr,no head,"\mathbb U^I"]&\\
\mathbb R^I\arrow[dr,no head,"\mathbb U^I"']&&S_*T^!(X)\\
&\mathbb Z^I\arrow[ur,"T^!(X)"']&
\end{tikzcd}
}
\margin{$T(X)\left\{\vphantom{\begin{matrix}a\\a\\a\end{matrix}}\right.$}
\margin{$\left.\vphantom{\begin{matrix}a\\a\\a\end{matrix}}\right\}\operatorname{Aut}((T_i)_{i\in I})$}

d'où
\diagram{
\begin{tikzcd}[row sep=large,column sep=small]
\operatorname{SRT}\deleted{\unsure{(X)}}/\mathbb Z^I
 \overset{\mathrm{def}}=\operatorname{SUT}(X)
 \arrow[d,"\unsure{\mathfrak S_I}"]\\
T^!(X)\times\mathbb U^I
\end{tikzcd}
}
\[
(*)\qquad1\longrightarrow\mathbb U^I\longrightarrow
\operatorname{SUT}(X)-\unsure{T|\Gamma}(X)-1
\]

De plus, chaque fois qu'on se donne un
marquage, i.e. un pt \deleted{\unsure{de}}
\[
\uline{s}=(s_i)\in
\prod_{i\in I=\pi_0(\partial X)}
\deleted{\unsure{\mathbb U_i}}T_i
\]
\margin{$T_i=$ comp. connexe (torique) de $\partial X$ d'indice $i$}

on trouve un plongement
\[
S_*T^{\unsure{!}}(X,\uline{s})\hookrightarrow\operatorname{SRT}(X)
\]
tel qu'on en déduise
\[
S_*T^{\unsure{!}}(X,\uline{s})
\mathbin{\wedge}_{\mathbb Z^I}\mathbb R^I
\simeq\operatorname{SRT}(X)
\]

\margin{1}
